\section{A worked decision-preservation analysis}\label{sec:C12:worked}
Consider a task with three obligations: modify a parser, preserve a legacy behavior, and produce a reviewable patch without publishing it. The raw history includes a failed test, a revised hypothesis, and a user correction that narrows the authorized output. A token-oriented summarizer may retain the patch plan and the latest test output while dropping the legacy requirement and the authorization correction as old discussion. The resulting input is shorter but misrepresents the permitted output scope and success requirements in the model-visible view.

Represent the validated state as $z=(G,O,A,E,U)$: goals $G$, obligations $O$, authority $A$, evidence references $E$, and unresolved hypotheses $U$. A proposed compaction operator $T$ should preserve $O$ and $A$ exactly, preserve enough of $E$ to justify current conclusions, and mark unresolved hypotheses as unresolved. This is an application-specific contract. It does not claim that these five fields are sufficient for every task; it states which distinctions this parser workflow requires.

Suppose the original and compacted policies induce action distributions $p$ and $q$ at the next decision. For a bounded immediate loss $0\leq\ell\leq1$, their expected losses differ by at most $\mathrm{TV}(p,q)$. If a validator can establish a bound on distribution change, that bound controls immediate risk under the stated loss. It does not automatically control the complete task, because a small action change can place the environment in a different region. Long-horizon guarantees require bounds on transition sensitivity or repeated coupling assumptions, which are rarely available for arbitrary tools.

For deterministic safety constraints, comparing distributions is unnecessary. Suppose every tool call passes the same monitor, whose allowed-action set is determined exclusively by the authoritative grant state and the current enforcement configuration. Holding those inputs fixed preserves the same syntactic boundary even when compaction changes the policy. Copying an old $A$ into a summary does not satisfy this condition after a grant is revoked; the monitor must consult current authority at dispatch, independently of the edited prose. Semantic correctness remains open: the agent can take an allowed but wrong action. This separation is valuable because it prevents summary quality from being responsible for all enforcement. Authority should survive in a typed control layer even when task evidence is compressed.

The corresponding experiment compares four views: full history, generic summary, typed-state summary, and typed state plus targeted raw retrieval. Keep the model, tools, and environment snapshot fixed. Randomize whether the legacy constraint appears early or late and whether the latest correction expands or narrows authority. Score patch correctness, publication behavior, retrieval use, and cost separately. This is a proposed experiment. No numerical outcomes are implied by the formulation.

A successful typed-state condition would support the mechanism that explicit obligations survive compression in this workload. It would not establish that a manually designed schema generalizes to unknown scientific or legal tasks. A failure would still be useful if it identifies a missing state field or an incorrect validation rule. The scientific objective is to expose which information transformations preserve behavior, rather than to award a generic summary the label of memory.

